\begin{abstract}
Large language models (LLMs) now define the state of the art in many reasoning, generation, and decision-support tasks, but their deployment cost, energy demand, latency, and hardware requirements remain substantial. Small language models (SLMs) provide an attractive alternative for edge deployment, privacy-sensitive settings, and task-specific optimisation, yet their capabilities are uneven across domains. This paper presents a systematic literature review of 54 studies on LLMs, SLMs, and SLM--LLM collaboration, with emphasis on hybrid, multi-agent, retrieval-centered, and monolithic architectures. The review is guided by four research questions addressing comparative performance, computational efficiency, orchestration mechanisms, and domain suitability. Using a PRISMA 2020-guided review design and a Web of Science-centered search process, we synthesize evidence on model choices, architectural patterns, datasets, metrics, baselines, latency, cost, energy, and reported limitations. The reanalysis of the corpus shows a strong shift toward hybrid and multi-agent designs together with a clear reporting asymmetry in which accuracy is common but cost and energy are rarely quantified. The findings indicate that SLMs are most competitive when they are paired with routing, retrieval, planning, or fusion mechanisms and when the problem space is narrow, structured, or domain-specific.

\keywords{systematic literature review \and large language models \and small language models \and hybrid architectures \and multi-agent systems \and edge-cloud inference}
\end{abstract}
